\chapter{Környezeti változók dokumentációja}
\label{ch:envvar}

A rendszer konfigurációja környezeti változókon keresztül történik, amelyeket a \texttt{.env} fájlok tartalmaznak. Külön fájl szolgál a Docker környezethez (\texttt{.env}), a fejlesztéshez (\texttt{.env.dev}) és az éles üzemhez (\texttt{.env.prod}).

\begin{table}[h!]
    \centering
    \small
    \renewcommand{\arraystretch}{1.3}
    \begin{tabularx}{\textwidth}{|l|X|X|}
        \hline
        \rowcolor[HTML]{EFEFEF} 
        \textbf{Változó} & \textbf{Alapértelmezett / Példa} & \textbf{Leírás} \\ \hline
        
        \multicolumn{3}{|l|}{\cellcolor[HTML]{F5F5F5}\textit{Docker (Alapkonfiguráció)}} \\ \hline
        \texttt{UID / GID} & \texttt{1000} & A gazdagép felhasználói és csoport azonosítója a fájljogosultságokhoz. \\ \hline
        
        \multicolumn{3}{|l|}{\cellcolor[HTML]{F5F5F5}\textit{Alkalmazás (Laravel)}} \\ \hline
        \texttt{APP\_NAME} & \texttt{SkillIssue} & Az alkalmazás megnevezése. \\ \hline
        \texttt{APP\_ENV} & \texttt{local / production} & A futtatási környezet típusa. \\ \hline
        \texttt{APP\_KEY} & \texttt{base64:...} & Az alkalmazás titkosító kulcsa (\texttt{make dev-key}). \\ \hline
        \texttt{APP\_URL} & \texttt{https://skillissue.hu} & Az alkalmazás elérési útvonala. \\ \hline
        \texttt{APP\_DEBUG} & \texttt{true / false} & Hibakeresési mód ki- és bekapcsolása. \\ \hline
        
        \multicolumn{3}{|l|}{\cellcolor[HTML]{F5F5F5}\textit{Adatbázis (MySQL)}} \\ \hline
        \texttt{DB\_CONNECTION} & \texttt{mysql} & Az adatbázis-kezelő típusa. \\ \hline
        \texttt{DB\_HOST} & \texttt{mysql} & Az adatbázis-szolgáltatás gépneve (Docker-en belül). \\ \hline
        \texttt{DB\_DATABASE} & \texttt{skillissue} & A használt adatbázis neve. \\ \hline
        \texttt{DB\_USERNAME} & \texttt{root / appuser} & Az adatbázis eléréséhez használt felhasználónév. \\ \hline
        \texttt{DB\_PASSWORD} & \texttt{********} & Az adatbázis jelszava. \\ \hline
        
        \multicolumn{3}{|l|}{\cellcolor[HTML]{F5F5F5}\textit{Játékmechanika}} \\ \hline
        \texttt{MAX\_ROUNDS} & \texttt{5} & Egy mérkőzés maximális köreinek száma. \\ \hline
        \texttt{SOLO\_MAX\_GUESS\_TIME} & \texttt{30} & Válaszadási idő másodpercben (gyakorló mód). \\ \hline
        \texttt{RANKED\_MAX\_GUESS\_TIME} & \texttt{30} & Válaszadási idő másodpercben (rangsorolt mód). \\ \hline
        \texttt{VITE\_*} & \texttt{...} & A frontend számára is elérhetővé tett változók (prefixelt). \\ \hline
        
        \multicolumn{3}{|l|}{\cellcolor[HTML]{F5F5F5}\textit{Biztonság és Kommunikáció}} \\ \hline
        \texttt{SERVICE\_TOKEN} & \texttt{abcdefg123456} & A belső szolgáltatások (API és WebSocket) közötti hitelesítéshez. \\ \hline
        \texttt{SANCTUM\_STATEFUL\_DOMAINS} & \texttt{skillissue.hu, ...} & Engedélyezett domainek a sütialapú hitelesítéshez. \\ \hline
    \end{tabularx}
    \caption{A projektben használt legfontosabb környezeti változók}
\end{table}

\noindent \textbf{Fontos:} Éles környezetben (\texttt{.env.prod}) a \texttt{SESSION\_ENCRYPT} értéke kötelezően \texttt{true}, a \texttt{LOG\_LEVEL} pedig \texttt{warning} szintre van állítva a biztonság és a teljesítmény érdekében.